%%%
%
% $Autor: Müller, Hinrichs $
% $Datum: 2026-06-09 14:18:35Z $
% $Pfad: Radar.tex
% $Version: 1 $
%
% Project: 
%
%%%


\chapter{Grundlagen der akustischen Ortungs- und Visualisierungstechnik}
\label{chap:Domain_Knowledge_Application}

Dieses Kapitel befasst sich mit den theoretischen und mathematischen Grundlagen der raumbezogenen Objektortung mittels Ultraschall. Es isoliert die zugrundeliegenden physikalischen Prinzipien der akustischen Distanzmessung von der konkreten Hardwarekomponente und mathematisiert die Transformation von Sensordaten in zweidimensionale Koordinatensysteme. Darüber hinaus werden systematische Messabweichungen, bedingt durch geometrische Signalaufweitungen, sowie industrielle Anwendungsfelder analysiert.

\section{Das akustische Echo-Laufzeitverfahren (Time-of-Flight)}
\label{sec:Time_of_Flight_Theorie}

Die Bestimmung geometrischer Abstände zwischen einer Sensorplattform und einer reflektierenden Grenzfläche basiert in der akustischen Ortungstechnik primär auf dem Prinzip des Echo-Laufzeitverfahrens, im internationalen Sprachgebrauch als \text{Time-of-Flight (ToF)} definiert \cite{Lerch:2012}. Ein elektroakustischer Wandler emittiert hierbei einen transienten, hochfrequenten Schallimpuls, welcher das umgebende Medium durchschreitet. Trifft diese Wellenfront auf ein Hindernis mit abweichender akustischer Impedanz, wird ein Teil der Schallenergie reflektiert und als Echo an den Wandler zurückgesendet.

Die raumzeitliche Distanz $d$ berechnet sich deterministisch aus der gemessenen Totalzeit $t_{total}$, welche das Signal für den Hin- und Rückweg benötigt, sowie der temperaturabhängigen Schallgeschwindigkeit $v_c$ innerhalb des gasförmigen Mediums \cite{Fraden:2010}:

\begin{equation}
	d = \frac{v_c \cdot t_{total}}{2}
\end{equation}

Für präzise mesSstechnische Anwendungen im Bereich des Maschinenbaus ist zu berücksichtigen, dass die Schallgeschwindigkeit in Luft ma\ss geblich von der thermodynamischen Temperatur $T$ in Grad Celsius ($^\circ$C) beeinflusst wird. Die Linearisierung der zugrundeliegenden Zustandsgleichung für ideale Gase ergibt für den praxisrelevanten Temperaturbereich die Approximation \cite{Lerch:2012}:

\begin{equation}
	v_c(T) \approx 331{,}5 \text{ m/s} + 0{,}6 \cdot T
\end{equation}

Systembedingte Latenzen, wie das Ausschwingverhalten der piezokeramischen Membranen (Blindzone) sowie elektronische Signalaufbereitungszeiten, limitieren dabei die minimale und maximale Detektionsreichweite auf软件ebene.

\section{Trigonometrische Koordinatentransformation für die Raumprojektion}
\label{sec:Koordinatentransformation}

Die sensorisch erfassten Rohdaten liegen nativ als diskrete Polarkoordinatenpaare vor, bestehend aus der skalaren Entfernung $r$ (äquivalent zu $d$) und dem mechanisch eingestellten Azimutwinkel $\theta$, welcher durch die Winkelstellung des Aktors vorgegeben wird. Für die Überführung dieser polaren Datensätze in ein computergestütztes, zweidimensionales Pixel-Raster (Grafik-Bildschirm) ist eine kontinuierliche Koordinatentransformation in ein kartesisches System erforderlich.

Die Transformation eines Punktes im euklidischen Raum erfolgt über die trigonometrischen Grundgleichungen:

\begin{equation}
	x_{radial} = r \cdot \cos(\theta)
\end{equation}
\begin{equation}
	y_{radial} = r \cdot \sin(\theta)
\end{equation}

Da computergestützte Visualisierungsumgebungen (wie Processing) ihren mathematischen Koordinatenursprung $(0,0)$ standardmä\ss ig in der linken oberen Bildschirmecke definieren und die vertikale Ordonnate ($Y$-Achse) nach unten gerichtet ist, muss eine translationale und spiegelnde Matrix-Transformation stattfinden. Soll der physische Drehpunkt des Radarsystems im unteren Zentrum des Bildschirms bei den Pixelkoordinaten $(x_0, y_0)$ liegen, lauten die modifizierten Abbildungsgleichungen für das Display-Raster \cite{Parisi:2014}:

\begin{equation}
	x_{screen} = x_0 + r \cdot \cos(\theta)
\end{equation}
\begin{equation}
	y_{screen} = y_0 - r \cdot \sin(\theta)
\end{equation}

Diese mathematische Anpassung stellt sicher, dass der dynamische Abtaststrahl eine exakt synchrone Bewegung zur realen mechanischen Ausrichtung des Sensors vollführt.

\section{Physikalische Grenzflächeneffekte und geometrische Divergenz}
\label{sec:Physikalische_Restriktionen}

In der industriellen Anwendung weisen akustische Ortungsverfahren spezifische, physikalisch bedingte Restriktionen auf, die bei der Interpretation der generierten Visualisierung berücksichtigt werden müssen. Im Gegensatz zu kohärenten optischen Systemen (Laser-Entfernungsmessung) breitet sich Ultraschall nicht als unendlich dünner Strahl, sondern in Form einer dreidimensionalen Keule aus. Diese geometrische Aufweitung wird als Divergenz bezeichnet und durch den material- und frequenzabhängigen Öffnungswinkel $\alpha$ charakterisiert \cite{Klemenz:2019}.

Diese Keulencharakteristik führt zu zwei signifikanten Artefakten in der Applikationsdomäne:

\begin{enumerate}
	\item \textbf{Geometrische Objektverbreiterung:} Trifft die Flanke des Schallkegels auf eine Objektkante, wird bereits vor Erreichen des tatsächlichen Objektzentrums ein Echo registriert. Das Objekt erscheint auf der Benutzeroberfläche winkelmä\ss ig künstlich verbreitert.
	\item \textbf{Gerichtete Reflexion (Spiegelung):} Trifft die Schallwelle in einem flachen Winkel auf eine glatte, harte Oberfläche, wird die Energie gemä\ss{} des Reflexionsgesetzes (Einfallswinkel gleich Ausfallswinkel) vollständig vom Sensor weggeleitet. Trotz der physischen Präsenz eines Hindernisses detektiert die Anwendung in diesem Sektor kein Echo (Signalverlust).
\end{enumerate}
	
	Zur Kompensation dieser physikalischen Unschärfen müssen softwareseitige Filteralgorithmen (z.\,B. Hysterese- oder Clusterfilter) implementiert werden, welche benachbarte Winkelsignale auf Kohärenz prüfen und Kantenübergänge mathematisch glätten.
	
	\section{Industrielle Relevanz und Applikationsdomänen}
	\label{sec:Industrielle_Relevanz}
	
	Das mechatronische Prinzip der Ultraschall-Raumerfassung ist aufgrund seiner Robustheit gegenüber optischen Störeinflüssen wie Staub, Nebel, Umgebungslicht oder transparenten Oberflächen ein fester Bestandteil moderner Automatisierungs- und Fahrzeugtechnik \cite{Fraden:2010}. 
	
	In der Automobilindustrie bildet das beschriebene ToF-Verfahren das technologische Fundament für Nahbereichs-Umfelderfassungen, primär im Rahmen von Parkdistanzkontrollen (PDC) und autonomen Einparkassistenten. In der industriellen Robotik und Intralogistik werden schallbasierte Ortungsapplikationen zur berührungslosen Hindernisvermeidung und Kollisionsschutzabsicherung bei Fahrerlosen Transportsystemen (FTS) eingesetzt, um die funktionale Sicherheit im gemeinsamen Arbeitsraum mit dem Menschen zu gewährleisten \cite{Klemenz:2019}.
